% §16.4 Adversarial Inputs
%
% Cross-reference resolution status:
%
% DEFINED here:
%   \label{sec:adversarial-inputs}            -- this section
%
% ALREADY DEFINED (chapter 4):
%   \Cref{ax:element-indep}                   -- chapter 4 §4.2 Axiom 1
%   \Cref{ax:block-rates-indep}               -- chapter 4 §4.3 Axiom 2
%   \Cref{sec:axiom-1}                        -- chapter 4 §4.2
%   \Cref{sec:axiom-2}                        -- chapter 4 §4.3
%
% ALREADY DEFINED (chapter 5):
%   \Cref{sec:bloom-from-scratch}             -- chapter 5 §5.2 Bloom construction
%
% INTERNAL refs (chapter 16):
%   \Cref{sec:varying-rates}                  -- §16.2
%   \Cref{sec:correlated-noise}               -- §16.3
%   \Cref{sec:where-framework-helps}          -- §16.5
%
% FORWARD refs (later chapters):
%   \Cref{ch:cipher-maps-trapdoor}            -- chapter 17 (Plan 0 stub)
%   \Cref{ch:encrypted-search}                -- chapter 18 (Plan 0 stub)
%
% §16.4 covers the third failure mode: the query stream is not iid;
% an adversary chooses the worst-case query. Both axioms can break;
% recap how. Give four diagnostic examples (Bloom-filter evasion,
% side-channel timing, adversarial ML, chosen-plaintext attacks).
% State the right tools (worst-case bounds, cryptographic hash
% families, oblivious primitives). Close with the Bloom-evasion
% worked example quantifying the iid-vs-adversarial gap.

\section{Adversarial Inputs}
\label{sec:adversarial-inputs}

The third failure mode breaks both axioms simultaneously, but in a
distinctive way. Where \Cref{sec:varying-rates} watched the channel
matrix drift through nature and \Cref{sec:correlated-noise} watched
queries carry shared noise sources, the present section watches
queries chosen with intent. An adversary observes the system's
responses and selects subsequent queries to expose the worst-case
behavior. The whole probabilistic setup, which presumed a query
distribution that the analyst could pin down in advance, gives way
to a worst-case query strategy that depends on the adversary's
goals and knowledge.

\subsection*{The phenomenon}

A query stream is adversarial when the choice of the $t$-th query
depends on the responses to queries $1, \ldots, t - 1$ and is
selected to maximize some attacker objective such as the empirical
false-positive rate, a side-channel signal, or a successful
collision. The defender's analysis used a probability distribution
over queries (often iid, sometimes a tractable Markov chain); the
attacker uses a strategy that exploits whatever structure the
defender's analysis assumed.

A simple way to picture the gap: a Bloom filter sized for
$\fprate = 0.01$ gives an iid attacker a $1\%$ chance of hitting a
false positive on each random query. An attacker who learns the
hash family can compute $h_1(x), \ldots, h_k(x)$ for candidate
inputs $x$ and select inputs whose hash images land at already-set
bit positions. The probability of a false positive per attempt
jumps from $0.01$ to close to $1$, conditional on a brief offline
search. The Bernoulli analysis described a rate of $0.01$ per
query; the adversarial reality is a rate near $1$ per crafted
query.

A more involved version of the same phenomenon is the
gradient-descent attacker. For a neural-network classifier whose
training error rate is some $\fprate$, an attacker who can query
the model and observe responses can use gradients to construct
inputs near decision boundaries. The empirical attack rate against
state-of-the-practice image classifiers is essentially $100\%$
with quite small perturbations; the iid analysis from labeled
benchmark data is a description of natural inputs, not of
adversarial inputs.

\subsection*{What's violated}

The adversarial setup fails both axioms, but the failures stack in
a specific way. \Cref{ax:element-indep} (Axiom 1) presumes that
the inputs $x_1, x_2, \ldots$ are conditionally independent given
the true labels; the attacker's $t$-th query is by construction a
deterministic function of the prior responses, so the inputs are
not independent of one another. The conditional-independence claim
of \Cref{sec:axiom-1} fails not through any property of the
channel but through the input distribution itself: a sequence of
queries that an adversary can adapt is not an iid draw.

\Cref{ax:block-rates-indep} (Axiom 2) is more subtly violated. The
per-block rates in the channel matrix are typically defined
relative to some input distribution. An adversary who selects
inputs at the tail of that distribution effectively replaces the
expected-case per-block rate with the worst-case per-block rate
over inputs in the relevant block. The channel still has a stable
matrix, but the matrix entries that the analyst computed for
``typical'' queries no longer describe the attacker's queries.
This is conceptually similar to \Cref{sec:varying-rates}'s
time-varying rates, except that the variation tracks adversarial
selection rather than time.

The cleanest way to describe the joint failure is that the
worst-case attacker effectively converts the analyst's
expectation-based statements about $\fprate$ into supremum-based
statements: where the analyst wrote $\mathbb{E}_x[\fprate(x)]$ for
an average over a presumed distribution, the attacker reads
$\sup_x \fprate(x)$ as the rate that actually obtains. The two
quantities can differ by orders of magnitude.

\subsection*{Concrete examples}

\textbf{Bloom filter evasion.} An attacker who knows the hash
family $h_1, \ldots, h_k$ used by a defender's Bloom filter $B_A$
can compute, for any candidate input $y$, the set of positions
$h_1(y), \ldots, h_k(y)$ that $B_A$ would read. If those positions
happen to all be set in $B_A$ (because other items in $A$
collectively set those bits), the filter returns $\top$ on $y$,
even though $y \notin A$. The attacker enumerates candidate inputs
until one returns $\top$. Under iid sampling from the universe,
finding such a $y$ takes $1 / \fprate$ tries on average. With
hash-family knowledge and a structured search, the attacker can
find a $y$ in a constant number of attempts by directly targeting
positions known to be set. The gap between $1 / \fprate$ and
$O(1)$ is the cost of unbounded adversarial knowledge.

\textbf{Side-channel timing.} A defender's filter may answer
queries faster or slower depending on which bits are read first.
Sequential bit-reads with early termination on a clear bit (the
filter can return $\bot$ as soon as one of the $k$ bits is found
unset) leak information about which bit positions are set: an
unusually slow response corresponds to all $k$ bits being checked,
i.e., a $\top$ response and therefore a candidate false-positive
target. An attacker who measures response time can localize the
attack region in a far smaller candidate space. The Bernoulli
framework's analysis described the bit-pattern probabilities but
not the access pattern; timing information sits outside the
framework's scope and licenses queries that the framework's
distribution does not account for.

\textbf{Adversarial ML inputs.} A classifier trained on a natural
distribution achieves a benchmark $\fprate$ on held-out data. An
attacker who can query the classifier finds inputs near decision
boundaries by following the classifier's confidence gradient (or
by using a substitute model as a proxy). Adversarial examples have
been a subject of sustained research; the canonical observation
from this literature is that the empirical worst-case error rate
under gradient attacks is essentially $100\%$ for many natural
classifiers, regardless of the benchmark $\fprate$ value. The
benchmark statement and the worst-case statement describe
different distributions of inputs; conflating them is the failure
mode.

\textbf{Cryptographic chosen-plaintext attacks.} In the encrypted
search context that \Cref{ch:encrypted-search} (chapter 18 \S18.3)
takes up, an attacker who can choose plaintexts and observe the
ciphertext or the index responses gets a much stronger position
than an attacker restricted to the natural distribution of
documents. The Bernoulli analysis for an encrypted index typically
assumed the document distribution; chosen-plaintext attacks let the
attacker pick the worst-case document, breaking the same
expectation-versus-supremum gap described above. Chapter 18
catalogues the standard cryptographic defenses (semantic security,
indistinguishability under chosen-plaintext attack); the present
section flags the structural form of the attack rather than the
cryptographic remediation.

\subsection*{The right tools}

When iid is not the right model, the analyst's two replacements
are worst-case bounds and cryptographic primitives.

\emph{Worst-case bounds.} Drop the expectation and bound the
maximum. For a Bloom filter, the worst-case $\fprate$ over inputs
is determined by the bit pattern of $B_A$ and is bounded above by
the fraction of bit-pattern subsets of size $k$ that lie inside
the set-bit positions. For a classifier, the worst-case loss over
a perturbation ball is bounded above by the Lipschitz constant of
the model times the ball radius. These bounds are typically
loose by large factors compared to natural-distribution
expectations, but they are the right object for adversarial
analysis: they describe what the attacker actually controls.

\emph{Cryptographic hash families.} A standard cryptographic move
is to make the hash family secret. The Bloom filter's hash family
$h_1, \ldots, h_k$ is generated using a secret salt
$\mathsf{salt}$, and the filter publishes only the bit pattern;
the salt remains with the defender. An attacker who tries to
construct false positives must first guess the salt, which under
standard hash-family assumptions requires $2^{|\mathsf{salt}|}$
work in expectation. The structural attack collapses to a
brute-force key search; the false-positive rate under salt-secret
operation matches the iid analysis, provided the salt has enough
entropy. This is the cryptographic core of cipher maps; chapter
17 picks up the details.

\emph{Privacy-preserving primitives.} For settings where the
attacker can observe side channels (timing, memory access
patterns, control flow), the right machinery is oblivious
computing: data-oblivious algorithms whose access patterns do not
depend on the data they process. The classical primitives are
oblivious RAM and constant-time cryptographic operations.
\Cref{ch:cipher-maps-trapdoor} (chapter 17 \S17.3) details
oblivious primitives as the bridge between Bernoulli analysis and
adversarial environments; the Bernoulli framework continues to
apply if and only if the side channels are closed.

\subsection*{Worked example: Bloom filter evasion}

Take a Bloom filter $B_A$ sized for $\fprate = 0.01$ at $k = 7$,
$m / n = 9.585$, storing $n = 1000$ items. The bit pattern has a
fraction $1 - e^{-kn/m} = 1 - e^{-0.73} \approx 0.518$ of bits set
at design fill. The iid false-positive rate of $0.01$ means that a
random non-member input has a $1\%$ chance of all $k$ hashed
positions being set, that is, $0.518^7 \approx 0.013$, which
matches the design target within rounding.

\emph{The iid attacker} (queries drawn uniformly at random)
expects to wait $1 / \fprate = 100$ queries to see one false
positive. The expectation has a geometric distribution with mean
$100$ and variance just under $10^4$; the iid attacker is a
diffuse, undirected adversary.

\emph{The structural attacker} (hash family known, no salt) can
do much better. The attacker knows $h_1, \ldots, h_k$ and can
read the bit pattern of $B_A$ (assume the filter is published; the
analyst's standard convention). For each candidate input $y$, the
attacker computes $h_1(y), \ldots, h_k(y)$ and checks whether all
$k$ positions are set in $B_A$. The attacker enumerates candidates
in any natural order until one satisfies the condition. Because
roughly $51.8\%$ of bits are set, a random candidate has
probability $0.518^7 \approx 0.013$ of working, and the attacker
expects to enumerate about $80$ candidates before a hit. The
iid-versus-structural gap is small here because the structural
attacker did not actually exploit additional knowledge beyond
hash-family disclosure; the worst-case attacker can do much
better.

\emph{The targeted attacker} (hash family known, plus optimization
over candidate space) can engineer inputs $y$ whose hashes land at
specific set positions. If the attacker can solve a small system
of constraints to construct $y$ given target positions
$j_1, \ldots, j_k$ that are all set, the attacker finds false
positives in $O(1)$ work per attempt. The expected number of
attempts drops to a small constant, and the practical attack
proceeds at machine speed.

The gap between iid analysis and worst-case adversarial analysis,
quantified here, ranges from a small constant (for the structural
attacker who exploits only hash-family disclosure) to several
orders of magnitude (for the targeted attacker who optimizes
candidates). In any case where the security model includes an
attacker with knowledge of the hash family, the iid $\fprate$
value is not a meaningful security parameter; the cryptographic
salt-secret construction or an oblivious primitive must be
invoked. \Cref{sec:where-framework-helps} returns to the question
of which framework constructs continue to give traction once these
defenses are in place.
